\section{Per-modality profiles}
\label{sec:s-profiles}
% Text and tables moved from the main text (Section 4.5) on 17 Nov 2026 to keep it at eight pages after the
% pretrained-model results were added; sec/main_profiles.tex holds the unchanged main-text tables.
Figure~\ref{fig:atlas-all} shows where in the image each rung is close or far, and Tables~\ref{tab:brain} to~\ref{tab:fundus} give the profiles of each modality at each set's original window (16\,mm cubes in the brain, 16 pixels in faces, retina and iris).
In the brain, the dataset average is closer to the target than the person's own opposite hemisphere under \ac{MSE} (0.413 against 0.717), \ac{SSIM} and \ac{lDDT}-like, and only \ac{TM}-like reverses this order\statuspre.
Wrong detail sits at a random person under every metric, and noise falls short of a random person under \ac{MSE} (0.921) but far beyond it under \ac{lDDT}-like (2.073)\statuspre.
In IXI, \acs{LPIPS} places the softened opposite hemisphere beyond the unchanged one (AlexNet 1.25 against 0.76, VGG 0.91 against 0.88) and passes T2 with both networks, which no error or similarity metric does in any set, while persistence gives both 0.80 and fails T2\statuspre.
In faces, a second photograph of the same person is further than the dataset average under \ac{MSE} (0.640 against 0.499 in CelebA), because pose and lighting change between photographs, while ArcFace places it closest (0.354)\statusfinal.
In the retina, the second visit of the same eye is at 0.49 in GRAPE, while the fellow eye is further than the dataset average (1.33 in GRAPE, 1.47 in PAPILA) and the mean of the 16 nearest people falls in between (0.86, 0.94)\statuspre.
The iris shows the same order with a larger gap: the same eye captured again is at 0.41, 0.57 and 0.18 in UTIRIS \ac{NIR}, UTIRIS \ac{VIS} and UPOL, the other eye of the same person is at 1.12 to 1.29, beyond the dataset average, and the nearest person in between\statuspre.
Across modalities, the opposite side of the same body is therefore not an intermediate rung between the same person and a random person: under squared error, it is usually farther from the population average, because mirrored anatomy differs in detail even within one person.
Only faces break this pattern, where the mirrored half is as close as the dataset average (0.99 on the $Q$ scale) and closer than a second photograph.
Relatives behave like random people under squared error: a sibling is at 1.84 times the dataset-average error in QTIM and 1.73 in IMDb-WIKI, close to the random-person level of about 2 (Supplementary Figure~S1)\statuspre.
% Sources: atlas_provenance.json (sibling Q, QTIM and IMDb-WIKI rows).
% Sources: tab/fundus_ratio.tex, tab/iris_ratio.tex, tab/dataset_card.tex.
% Sources: fig_metric_summary_f8_R.json; common_all.json (QTIM, IXI); face_ladder_metrics.json; rq4_window_size.json (size curve, fig/rq4_size_plotted.json).

\begin{table}[htbp]
\centering
\caption{\textbf{Brain (QTIM).}
$R$ of Eq.~2 of the main text anchored on the second scan (0 = second scan, 1 = average distance of 16 random people); opposite side = opposite hemisphere; mean and 95\% interval from resampling families and positions.
Bold: closest fill; underline: second closest.
Under \acs{MSE}, \acs{SSIM} and \acs{lDDT}-like the dataset average is closer than the person's own opposite hemisphere, and wrong detail sits at a random person.}
\label{tab:brain}
\scriptsize
\setlength{\tabcolsep}{2pt}
\begin{tabular}{@{}lcccc@{}}
\toprule
Metric & Dataset average & Opposite side & Wrong detail & Noise \\
\midrule
MSE & \shortstack{\textbf{0.413}\statuspre\\{\scriptsize [0.367, 0.443]}} & \shortstack{\underline{0.717}\statuspre\\{\scriptsize [0.622, 0.792]}} & \shortstack{1.006\statuspre\\{\scriptsize [0.992, 1.020]}} & \shortstack{0.921\statuspre\\{\scriptsize [0.893, 0.944]}} \\
SSIM & \shortstack{\textbf{0.738}\statuspre\\{\scriptsize [0.680, 0.784]}} & \shortstack{\underline{0.823}\statuspre\\{\scriptsize [0.755, 0.873]}} & \shortstack{1.023\statuspre\\{\scriptsize [1.014, 1.036]}} & \shortstack{1.132\statuspre\\{\scriptsize [1.088, 1.194]}} \\
\shortstack[l]{Residual\\correlation} & \textit{undefined} & \shortstack{\textbf{0.752}\statuspre\\{\scriptsize [0.688, 0.808]}} & \shortstack{\underline{1.003}\statuspre\\{\scriptsize [0.996, 1.009]}} & \shortstack{\underline{1.003}\statuspre\\{\scriptsize [1.000, 1.006]}} \\
TM-like & \shortstack{\underline{0.723}\statuspre\\{\scriptsize [0.642, 0.778]}} & \shortstack{\textbf{0.667}\statuspre\\{\scriptsize [0.530, 0.761]}} & \shortstack{0.997\statuspre\\{\scriptsize [0.987, 1.007]}} & \shortstack{1.075\statuspre\\{\scriptsize [1.041, 1.121]}} \\
\shortstack[l]{lDDT-like,\\repeat tol.} & \shortstack{\textbf{0.355}\statuspre\\{\scriptsize [0.254, 0.430]}} & \shortstack{\underline{0.890}\statuspre\\{\scriptsize [0.843, 0.925]}} & \shortstack{1.050\statuspre\\{\scriptsize [1.030, 1.080]}} & \shortstack{2.073\statuspre\\{\scriptsize [1.924, 2.273]}} \\
\bottomrule

\end{tabular}
\end{table}
% Source: results/donor_ladder_brain_metrics.json.

\begin{table}[htbp]
\centering
\caption{\textbf{Faces.}
$R$ of Eq.~2 of the main text anchored on the target itself; nearest person = the single unrelated person with the closest context ring; ArcFace compares whole images.
Bold: closest fill; underline: second closest.
A local nearest person is closest under \acs{MSE}, yet ArcFace places it almost at a random person (0.98), while the person's own second photograph is closest under ArcFace.}
\label{tab:face}
\scriptsize
\setlength{\tabcolsep}{2pt}
\begin{tabular}{@{}llcccc@{}}
\toprule
Data & Metric & 2nd photo & Opp. half & Average & Nearest person \\
\midrule
CelebA & MSE & 0.640\statusfinal & \underline{0.475}\statusfinal & 0.499\statusfinal & \textbf{0.399}\statusfinal \\
CelebA & SSIM & 0.828\statusfinal & \textbf{0.693}\statusfinal & 0.844\statusfinal & \underline{0.726}\statusfinal \\
CelebA & ArcFace & \textbf{0.354}\statusfinal & \textit{unavailable} & \textit{unavailable} & \underline{0.983}\statusfinal \\
IMDb-WIKI & MSE & 0.690\statusfinal & 0.529\statusfinal & \underline{0.498}\statusfinal & \textbf{0.397}\statusfinal \\
IMDb-WIKI & SSIM & 0.856\statusfinal & \textbf{0.739}\statusfinal & 0.859\statusfinal & \underline{0.743}\statusfinal \\
IMDb-WIKI & ArcFace & \textbf{0.396}\statusfinal & \textit{unavailable} & \textit{unavailable} & \underline{0.984}\statusfinal \\
\bottomrule

\end{tabular}
\end{table}
% Source: results/face_ladder_metrics.json.

\begin{table}[htbp]
\centering
\caption{\textbf{Retina and iris} ($Q$: 1 = dataset average, about 2 = random person).
Retina: 16 nearest = mean of the 16 nearest people; PAPILA has no second visit.
Iris: nearest person = one unrelated eye with the closest context.
Bold: closest fill; underline: second closest.
The person's own other eye is farther than the dataset average in every set, and only the same eye seen again comes clearly closer.}
\label{tab:fundus}
\label{tab:iris}
\scriptsize
\setlength{\tabcolsep}{2pt}
\begin{tabular}{@{}lcccc@{}}
\toprule
Retina & Second visit & Fellow eye & 16 nearest & Random person \\
\midrule
GRAPE P8 & \shortstack{\textbf{0.489}\statuspre\\{\scriptsize [0.404, 0.585]}} & \shortstack{1.330\statuspre\\{\scriptsize [1.224, 1.437]}} & \shortstack{\underline{0.856}\statuspre\\{\scriptsize [0.832, 0.880]}} & \shortstack{2.012\statuspre\\{\scriptsize [1.897, 2.140]}} \\
PAPILA P9 & \textit{unavailable} & \shortstack{\underline{1.467}\statuspre\\{\scriptsize [1.422, 1.512]}} & \shortstack{\textbf{0.935}\statuspre\\{\scriptsize [0.914, 0.955]}} & \shortstack{1.992\statuspre\\{\scriptsize [1.941, 2.049]}} \\
\bottomrule

\end{tabular}

\smallskip
\begin{tabular}{@{}lcccc@{}}
\toprule
Iris & Same eye & Other eye & Nearest person & Random person \\
\midrule
UTIRIS NIR & \shortstack{\textbf{0.407}\statuspre\\{\scriptsize [0.345, 0.475]}} & \shortstack{1.256\statuspre\\{\scriptsize [1.123, 1.395]}} & \shortstack{\underline{0.599}\statuspre\\{\scriptsize [0.540, 0.660]}} & \shortstack{2.001\statuspre\\{\scriptsize [1.816, 2.207]}} \\
UTIRIS VIS & \shortstack{\textbf{0.569}\statuspre\\{\scriptsize [0.503, 0.640]}} & \shortstack{1.293\statuspre\\{\scriptsize [1.211, 1.375]}} & \shortstack{\underline{0.861}\statuspre\\{\scriptsize [0.826, 0.898]}} & \shortstack{2.018\statuspre\\{\scriptsize [1.884, 2.176]}} \\
UPOL & \shortstack{\textbf{0.177}\statuspre\\{\scriptsize [0.154, 0.200]}} & \shortstack{1.118\statuspre\\{\scriptsize [1.037, 1.198]}} & \shortstack{\underline{0.841}\statuspre\\{\scriptsize [0.798, 0.882]}} & \shortstack{2.138\statuspre\\{\scriptsize [2.007, 2.281]}} \\
\bottomrule

\end{tabular}
\end{table}
% Sources: results/donor_ladder_2d_facefmt.json; results/iris_ladder_metrics.json.
